\subsection{Internal ring localization}
\label{sec:annex-p10-internal-ring-fractions}

\begin{annexitem}{annex:p10-internal-ring-fractions}{thm:p10-internal-ring-fractions}{Fractions without choice of representatives}
\begin{proof}
All equations and existential arguments below are internal to the topos.
They are intuitionistically valid.

\emph{The equivalence relation.}
Reflexivity uses witness \(1\in S\), and symmetry uses the same
witness with the sign of the difference reversed.
Suppose \(u(at-bs)=0\) and \(v(br-ct)=0\), with \(u,v\in S\).
Multiplying the first equation by \(vr\), the second by \(us\),
and adding gives
\[
 uv\,t(ar-cs)=0.
\]
The element \(uvt\) belongs to \(S\), so it witnesses
\((a,s)\sim(c,r)\).
Existential elimination supplies this argument from the two
existential premises; no witness-selection function is required.
The relation is consequently an internal equivalence relation.
Its effective quotient exists by
\cref{thm:p10-lawvere-sheaves-effective-quotients}.
Effectivity ensures that equality of two represented fractions is
exactly the displayed existential relation.

\emph{Descent of the operations.}
Suppose \((a,s)\sim(b,t)\) with witness \(u\), and
\((c,r)\sim(d,q)\) with witness \(v\).
For addition the cross-multiplied difference is
\[
 (ar+cs)tq-(bq+dt)sr
 =rq(at-bs)+st(cq-dr).
\]
Multiplication by \(uv\in S\) makes it zero.
For multiplication the corresponding difference is
\[
 ac\,tq-bd\,sr=cq(at-bs)+bs(cq-dr),
\]
again annihilated by \(uv\).
Negation changes the sign of the cross-multiplied difference,
so the original witness suffices.
The proposed operations therefore respect the equivalence relation.
Products of the quotient epimorphism are epimorphisms, and the
effective-quotient property descends these compatible operations
to arrows on \(S^{-1}A\).

For completeness, scaling numerator and denominator by any
\(v\in S\) does not change the class: the cross-multiplied
difference between \((a,s)\) and \((av,sv)\) is zero, so
witness \(1\) suffices. This observation also permits comparison
of expressions with different common denominators.

\emph{The ring laws.}
Addition is commutative because its numerator is \(at+bs\)
and its denominator is \(st\).
Both associations of the sum of \(a/s,b/t,c/r\) have representative
\[
 (atr+bsr+cst,\;str).
\]
Multiplication is commutative, and both associations of a triple
product have representative \((abc,str)\).
The additive zero and multiplicative unit laws follow by substituting
\(0/1\) and \(1/1\).
The sum of \(a/s\) and its negative is \(0/s^2=0/1\).
For distributivity, \(a/s\) times \(b/t+c/r\) has representative
\((a(br+ct),str)\).
The sum of the two separate products has representative
\((s\,a(br+ct),s^2tr)\), which is the same fraction by scaling.
Thus the ring identities hold on represented elements and hence
hold on the quotient: equality of arrows can be checked after
precomposition with the appropriate epimorphism.
The map \(i(a)=a/1\) preserves all ring operations.
For \(s\in S\), the fraction \(1/s\) is an inverse of \(i(s)\),
since \(s/s=1/1\).

\emph{The universal property.}
Let \(h:A\to B\) invert \(S\).
The inverse of an invertible element is unique, so inversion
defines an arrow on the unit subobject of \(B\).
If \(u(at-bs)=0\) witnesses equality of fractions, application
of \(h\) and cancellation of the unit \(h(u)\) gives
\(h(a)h(t)=h(b)h(s)\).
Multiplying by the two denominator inverses shows
\[
 h(a)h(s)^{-1}=h(b)h(t)^{-1}.
\]
The arrow on \(A\times S\) with this value therefore descends
uniquely through the quotient. The fraction operation formulas
show that its descended arrow \(\bar h\) is a ring homomorphism
and that \(\bar h\,i=h\).
Conversely every fraction equals \(i(a)i(s)^{-1}\).
Any ring homomorphism extending \(h\) must send that expression
to \(h(a)h(s)^{-1}\), because it preserves inverses of units.
The quotient epimorphism then gives uniqueness.

\emph{Change of topos.}
An inverse-image functor of a geometric morphism preserves finite
limits and coequalizers. It therefore preserves regular-epi--mono
image factorizations and effective quotients.
The fraction relation is constructed using finite products,
ring-operation arrows, equality, and the image expressing
\(\exists u\in S\). Its inverse image is exactly the fraction
relation for the transported ring and multiplicative subobject.
The transported quotient and operations are accordingly the
localization in the target topos. The resulting isomorphism is
the canonical one determined by the universal property.
No preservation of an arbitrary power object is used.
\end{proof}
\end{annexitem}

\begin{annexitem}{annex:p10-internal-prime-locality}{thm:p10-internal-prime-locality}{Units and prime predicates}
\begin{proof}
First the numerator criterion descends to fractions.
If \(u(at-bs)=0\) and \(a\in S\), then \(uat=ubs\).
The left side belongs to \(S\), so the right side does too.
The multiplication equivalence for \(S\) implies \(b\in S\).
The reverse implication follows by symmetry.

If \(a\in S\), the fraction \(s/a\) is defined and inverts
\(a/s\). Conversely, if \(a/s\) is a unit, take an inverse
locally in represented form \(b/t\). Such local representation
suffices because the quotient map is an epimorphism.
Equality \(ab/(st)=1/1\) means that locally some \(u\in S\)
satisfies \(u(ab-st)=0\). Hence
\[
 uab=ust\in S.
\]
The multiplication equivalence implies \(a\in S\).
This proves the unit criterion internally, without choosing
global representatives.

The fraction \(1/1\) is a unit. The fraction \(0/1\) is not
a unit, since the criterion would give \(0\in S\).
Equivalently, equality \(0/1=1/1\) would have a witness
\(u\in S\) with \(u=0\), contradicting \(0\notin S\).
For fractions \(a/s,b/t\), their product is a unit precisely
when \(ab\in S\), and this is equivalent to both \(a\in S\)
and \(b\in S\). If their sum is a unit, then \(at+bs\in S\).
The addition condition gives \(at\in S\) or \(bs\in S\);
the multiplication equivalence gives \(a\in S\) or \(b\in S\).
The unit criterion thus verifies all four prime conditions
for the unit subobject of the localized ring. The disjunction
remains an internal disjunction and is not replaced by a
global decision about which alternative holds.

Finally let \(h:A\to B\) have local target, and let \(S\)
be the inverse image of the units of \(B\).
The unit subobject of a local ring satisfies the four prime
conditions by definition. Pulling those conditions back along
the ring homomorphism proves them for \(S\).
In particular \(h\) inverts \(S\), so the preceding universal
property gives the unique factorization through \(S^{-1}A\).
For a represented fraction its image is \(h(a)h(s)^{-1}\).
Because \(h(s)\) is a unit, this image is a unit if and only
if \(h(a)\) is a unit, equivalently \(a\in S\), equivalently
\(a/s\) is a unit. The induced map therefore reflects units.
Every ring homomorphism preserves units, so both claims hold.
\end{proof}
\end{annexitem}
